\section{Finite Exact Search and Original-Instance Certificates}\label{sec:branch54}
\subsection{Price bounds, feasible policies, and discrete coverage}
For a restricted family $(M,F)$, every evaluated price yields an upper bound $U_{M,F}(\lambda)$. Recovering the price-maximizing book at $B$ gives a feasible lower value. Neither the priced mean nor a solver's completion status is a feasible-policy witness. Maintain a global incumbent $\underline V$ and, for each live leaf, its best completed price bound. The gross terminal-reward bound $\sum_j\pi_jf_j(1)$, dropping nonnegative service costs and opening charges, initializes a leaf before its first price evaluation.

The reference search evaluates a bounded number of prices at each node. The target range of a maximizing book gives a valid subgradient direction for the convex price bound. A bracket starts at $-\max_j\gamma_j$ and $\max_jr_j$; candidate prices also include the exact supporting price of the recovered fixed-book policy. Bisection supplies a distinct rational price when a candidate repeats. Bracketing is a search heuristic, not a premise for validity: every completed price produces a valid upper bound independently of how it was chosen.

If a node bound is at most $\underline V+\varepsilon$, it need not be split. Otherwise select an undecided candidate and create the two children in which it is respectively required and forbidden. Both inherit the parent's bound until a sharper completed calculation is available. Mandatory and forbidden vertices are enforced by the recurrence in Section~\ref{sec:price54}. A family with one possible book is evaluated at that book's exact supporting price. Every feasible original book remains in exactly one leaf family.

\begin{theorem}[Grid-free finite exact algorithm]\label{thm:branch54}
Under the rational quadratic assumptions of Proposition~\ref{prop:oracle54}, the preceding algorithm returns, after every completed node or price calculation, an original feasible policy and a rational interval
\begin{equation}\label{eq:anytime54}
 \underline V\leq V^*(B)\leq\overline V,
 \qquad
 \overline V=\max\{\underline V,\ U_{M,F}: (M,F)\text{ is a current leaf}\}.
\end{equation}
An empty leaf is omitted. With no resource limit and zero requested width, finite binary branching and exact singleton-family evaluation terminate at $\underline V=\overline V$. The full binary cover has at most $2^{N+1}-1$ nodes. A node may use multiple prices, so total nodes and completed price calls are different counts. Each completed price oracle has the polynomial complexity of Proposition~\ref{prop:oracle54}. At positive tolerance, or under interruption, the actual width $\overline V-\underline V$ is a valid instance-dependent additive error bound; it has no mandatory inverse-tolerance resource grid.
\end{theorem}
\begin{proof}
The initial incumbent is an original feasible singleton at the minimum command. Every subsequent incumbent is recovered at $B$ and respects the same constraints. Theorem~\ref{thm:price54} bounds every book in a leaf. Requiring or forbidding a previously undecided candidate partitions its family exactly, so inherited bounds remain valid on both children. A node with no allowed path is empty; all nonempty leaf families have feasible target allocations by Proposition~\ref{prop:canonical52}. These facts prove \eqref{eq:anytime54}, including leaves not yet processed or calculations interrupted before completion.

Along a root-to-leaf chain, a candidate is decided at most once. After at most $N$ splits the family is empty or fixes one book. Concave fixed-book allocation has a supporting price at $B$ (at a boundary promise, a one-sided supergradient extended to the feasible interval suffices, and the quadratic model has finite marginal bounds), and its scalar maximizers attain the price bound at that same promise. Consequently the restricted priced optimum at that price equals the recovered value. A fixed finite number of price attempts per nonterminal node followed by an undecided split cannot prevent this finite termination. The elementary depth-$N$ binary-tree count supplies the worst-case bound. Rational formulas, finite supports, and original-space policies establish exact arithmetic throughout.
\end{proof}

The worst-case exponential bound is deliberate. Theorem~\ref{thm:hard54} proves NP-completeness of the unrestricted decision problem. The present algorithm supplies a finite exact method with a polynomial node relaxation and a data-dependent certified stopping rule; it does not assert a polynomial worst-case exact running time. The exact common-cap and common-prefix theorems give polynomial subcases. For either reward model with the appropriate arc kernels, the uniform-resource scheme in Section~\ref{sec:uniform54} retains a separate worst-case polynomial additive guarantee. A valid combined interval can take the maximum of the feasible lower values and the minimum of the price-tree and resource-grid upper bounds. No assumption of strong duality across books is required.

\subsection{Independent verification and proof economics}
A price certificate contains the original model, its catalog and charges, a feasible lottery policy, and the complete binary coverage tree. A priced leaf records its rational price and upper bound; a still-unprocessed leaf may retain the universal bound. The independent checker reconstructs arc supports by direct endpoint and stationary-point maximization, then uses a backward at-most-budget recurrence rather than the optimizer's forward exact-length recurrence. It checks every branch, every mandatory or forbidden restriction, and every leaf bound. A parent price bound may be looser than a child's recomputed bound; the required comparison is dominance, not equality.

The lower witness retains nonnegative service, lottery masses, command support, every positive-probability realization ceiling, all expected caps, and the exact weighted promise. The checker subtracts the original selected charges and recomputes the global interval. A separately supplied instance digest can bind the proof to the requested input; a self-reported digest alone is not authentication. The new checker raises explicit errors and does not rely on assertions being enabled. The inherited grid checker continues to require ordinary Python execution, as documented in its retained source.

The proof record stores $O(k+s+N+P)$ rational fields and tree entries up to the history-level lottery representation, where $P$ is the number of retained nodes. It need not serialize the scalar-resource tables. Verification recomputes polynomial price paths at the leaf prices; it is not free, and a large branching tree can still be expensive. Input and price bit lengths contribute to actual byte counts. The study reports compressed and uncompressed size, checking time, peak process memory, and resource-limit outcomes separately from the command alphabet.

\subsection{Incumbent-based exclusion, including anchors}\label{sec:anchor54}
The uniform anchor-envelope rule retained in the electronic companion cannot delete a possible anchor. The priced oracle provides a complementary test.
\begin{proposition}[Forced-command price exclusion]\label{prop:fix54}
Let $v_0$ be the value of an original feasible incumbent. For any candidate $z$, including a possible anchor, let $U_z(\lambda)$ be the price-path upper bound with $z$ mandatory. If $U_z(\lambda)<v_0$, no optimal book contains $z$. All commands certified by such strict inequalities may be removed simultaneously without changing the optimal value. Under a weak inequality, simultaneous deletion is safe only with an independently retained incumbent avoiding all deleted commands.
\end{proposition}
\begin{proof}
Every book containing $z$ lies in its forced family and has value at most $U_z(\lambda)$. A strict comparison excludes it from optimality and also implies that the incumbent does not contain it. Therefore every optimal book and the incumbent avoid the union of strictly excluded candidates. The reduced problem retains an optimal policy and introduces no new one. For weak inequalities, any book lost by a deletion has value at most $v_0$, while an incumbent avoiding the entire deletion set attains $v_0$ in the retained problem. This proves that qualified version as well.
\end{proof}

This is the familiar incumbent-based Lagrangian reduction principle \citep{BeasleyChristofides1989,Fisher1981}, applied through an exact original-instance price oracle. Its benefit here is that the bound incorporates the accepted promise, command budget, and other selected alternatives. No uniform removal fraction follows. Reusing the arc supports for one price, evaluating every forced candidate costs $O(mN^3)$ additional arithmetic operations. The cheap $O(kN)$ non-anchor envelope remains a different option with an order-independent target-box guarantee. The screening study reports both costs and compares their exclusions with exact inclusion relevance rather than declaring one rule universally preferable.
